\section{Method}

\subsection{Column Vectorization}

Let $D$ be an input dataframe with $m$ rows and $d$ columns. Numeric columns are
converted to floating point values and missing values are imputed with the
column median. For a non-numeric column $c$, DataFrameSampler supports two
vectorization strategies.

If no helper columns are provided, each category is represented by its empirical
frequency. If helper columns are provided, each row receives a helper vector
$h_i \in \mathbb{R}^{k}$ constructed from the selected numeric helper columns.
An embedding function
\[
  f: \mathbb{R}^{k} \rightarrow \mathbb{R}
\]
maps each helper vector to a scalar used to discretize column $c$. In the
implementation, $f$ may be selected from several scikit-learn reducers,
including PCA, MDS, Kernel PCA, Isomap, locally linear embedding, and related
methods. Users may also supply custom objects implementing a scikit-learn-like
\texttt{fit\_transform}, \texttt{fit}/\texttt{transform}, or \texttt{transform}
interface.

\subsection{Discretization And Decoding}

After vectorization, each column is discretized into up to $b$ bins using a
one-dimensional binning procedure. Each row $x_i$ is encoded as an integer
latent vector
\[
  z_i \in \{0, \ldots, b-1\}^{d}.
\]
For every column and bin, the encoder stores the original observed values that
fell into that bin. Decoding a generated latent vector therefore consists of
sampling original values from the selected bins. If a generated bin is empty,
the implementation searches for the nearest non-empty bin in that column.

\subsection{Neighbour-Based Latent Sampling}

Let $Z = \{z_i\}_{i=1}^{m}$ be the encoded latent matrix. DataFrameSampler builds
a nearest-neighbour graph in this latent space. By default, only mutual
nearest-neighbour relationships are retained. If a row has no mutual neighbours,
a nearest-neighbour fallback is used to preserve sampleability.

To generate a new latent row, the sampler selects an anchor row $z_a$, a
neighbour $z_b$ of that anchor, and a neighbour $z_c$ of $z_b$. It then forms
\[
  \tilde{z} = z_a + \alpha (z_c - z_b),
\]
where $\alpha$ is sampled from a configured interpolation range. The generated
latent vector is rounded to integer bin indexes, clipped to valid ranges, and
decoded back to the original dataframe space.

\subsection{LLM-Assisted Configuration}

Choosing helper columns for categorical vectorization can require semantic
knowledge of column roles. The package therefore includes an optional
OpenAI-assisted mode. The helper sends a compact dataframe profile to an LLM,
including column names, dtypes, missing counts, unique counts, examples, and
basic numeric summaries. The model returns a structured configuration
containing candidate helper columns, sampled columns, an embedding method, and a
nearest-neighbour backend. The command-line flag \texttt{-A} fills only options
that the user did not explicitly specify.

\subsection{Sensitive-Value Replacement}

For workflows involving direct identifiers, the package can replace selected
sensitive columns before fitting the sampler. In this mode, OpenAI-generated
surrogate values are inserted into the original dataframe, exact and normalized
overlap with the source values is checked, and the sampler is then fit on the
surrogate dataframe. Generated samples are checked again against the original
source values. This procedure reduces exact identifier re-emission in selected
columns, but it is not a formal privacy guarantee.
